% =====================================================================
%  Documento "wrapper" per l'ANTEPRIMA delle bozze italiane dei
%  Capitoli 3 e 4. NON e' la tesi finale: serve solo a compilare in PDF
%  i frammenti di capitolo (che da soli non hanno un preambolo).
%
%  Compilazione:
%     pdflatex anteprima_bozze_it.tex      (due volte, per ToC e riferimenti)
%  oppure:
%     latexmk -pdf anteprima_bozze_it.tex
% =====================================================================
\documentclass[12pt,a4paper,oneside]{report}

\usepackage[utf8]{inputenc}
\usepackage[T1]{fontenc}
\usepackage[a4paper,margin=2.7cm]{geometry}
\usepackage{graphicx}
\usepackage{amsmath,amssymb}
\usepackage{booktabs}
\usepackage{array}
\usepackage{xcolor}
\usepackage{caption}
\usepackage{subcaption}
\usepackage{listings}
\usepackage{algorithm}
\usepackage{algpseudocode}
\usepackage{enumitem}
\usepackage{tabularx}
\usepackage{siunitx}
\usepackage[hidelinks]{hyperref}

\definecolor{codebg}{rgb}{0.97,0.97,0.97}
\definecolor{codekw}{rgb}{0.10,0.30,0.65}
\definecolor{codecm}{rgb}{0.30,0.55,0.30}
\lstset{
  basicstyle=\ttfamily\small,
  backgroundcolor=\color{codebg},
  keywordstyle=\color{codekw}\bfseries,
  commentstyle=\color{codecm}\itshape,
  stringstyle=\color{purple},
  numbers=left, numberstyle=\tiny\color{gray},
  breaklines=true, frame=single, framesep=4pt,
  showstringspaces=false, captionpos=b, language=Python,
}

\newcommand{\bimin}{BI_{\min}}
\newcommand{\bimax}{BI_{\max}}
\newcommand{\nthr}{N_{\mathrm{thr}}}
\renewcommand{\arraystretch}{1.15}

% Placeholder per il riferimento al Capitolo dei risultati (non ancora scritto),
% cosi' \ref{ch:risultati} non resta "??" in anteprima.
\newcommand{\chrisultati}{5}

\begin{document}

\tableofcontents
\clearpage

% =====================================================================
%  BOZZA - Capitolo 3 (italiano)
%  Il sistema ProSafe: protocolli di beaconing e meccanismi multihop
%
%  Richiede i pacchetti gia' presenti nel preambolo di main.tex:
%  amsmath, amssymb, siunitx, enumitem, booktabs.
%  Le macro seguenti sono ridefinite con \providecommand per sicurezza.
% =====================================================================
\providecommand{\bimin}{BI_{\min}}
\providecommand{\bimax}{BI_{\max}}
\providecommand{\nthr}{N_{\mathrm{thr}}}

\chapter{Il sistema ProSafe: protocolli di beaconing e meccanismi multihop}
\label{ch:sistema}

Questo capitolo descrive il sistema \emph{ProSafe} e i suoi protocolli di
beaconing. Si presentano l'architettura e i principi di progettazione, il modello
di comunicazione (struttura del beacon e accesso al canale), i tre scheduler
considerati --- il protocollo statico di riferimento (SBP), l'adattamento basato
sulla densit\`a (ADAB) e il protocollo context-aware (ACAB) --- e i due meccanismi
di consapevolezza multihop (\emph{append} e \emph{forwarded}).

% ---------------------------------------------------------------------
\section{Panoramica e principi di progettazione}
\label{sec:overview}

ProSafe (\emph{Proximity-based Safety}) \`e una rete di sicurezza basata su Wi-Fi
in cui delle \emph{boe intelligenti} (\emph{smart buoys}) trasmettono
periodicamente beacon broadcast con la propria posizione, cos\`i che le
imbarcazioni possano rilevare le persone in acqua. Il sistema lavora all'interno
di un dominio ad hoc IEEE~802.11 locale: gli allarmi sono generati e gestiti sul
posto, senza access point, router o connettivit\`a Internet.

La progettazione segue cinque principi:
\begin{itemize}[leftmargin=1.4em,nosep]
  \item \textbf{Assenza di infrastruttura}: nessun access point, router o rete
  cellulare. Il sistema deve funzionare in zone costiere dove non esiste alcuna
  infrastruttura.
  \item \textbf{Comunicazione broadcast}: il beaconing avviene a livello MAC con
  frame broadcast, senza i protocolli di livello IP.
  \item \textbf{Rilevazione passiva}: le unit\`a di bordo ascoltano senza
  trasmettere.
  \item \textbf{Hardware a basso costo}: si usano chipset Wi-Fi di largo consumo
  (smartphone, single-board computer).
  \item \textbf{Efficienza energetica}: le boe a batteria devono funzionare per
  ore con energia limitata, quindi la frequenza di trasmissione va gestita con
  attenzione.
\end{itemize}

\subsection{Componenti del sistema}
L'architettura prevede due tipi di dispositivo. La \emph{smart user buoy} \`e un
dispositivo galleggiante indossato o trasportato da chi \`e in acqua (nuotatori,
snorkeler, kayaker, sub). Integra un single-board computer con radio Wi-Fi (ad
esempio un Raspberry~Pi), un modulo GPS, un'unit\`a inerziale, una batteria e un
involucro stagno. La sua unica radio trasmette i propri beacon e ascolta quelli
delle altre boe, in modalit\`a ad hoc e senza alcuna associazione di rete. La
\emph{vessel safety unit}, installata a bordo delle imbarcazioni, lavora solo in
ascolto: quando riceve un beacon con potenza superiore a una soglia configurabile
genera un allarme acustico e visivo e pu\`o mostrare le boe rilevate su una mappa.
Poich\'e ogni beacon contiene anche la lista dei vicini recenti del mittente,
l'imbarcazione pu\`o accorgersi di boe che non ha sentito direttamente.

% ---------------------------------------------------------------------
\section{Modello di comunicazione}
\label{sec:commodel}

ProSafe esegue la scoperta di prossimit\`a a livello MAC IEEE~802.11, senza access
point, associazione o servizi di livello IP.

\subsection{Struttura del beacon}
\label{sec:beacon}
I beacon sono brevi frame broadcast inviati all'indirizzo di broadcast. Il payload
modellato nel simulatore (classe \texttt{Beacon}) contiene:
\begin{itemize}[leftmargin=1.4em,nosep]
  \item \textbf{sender id} --- identificatore del mittente ($16$ byte);
  \item \textbf{flag di mobilit\`a} --- indica se il nodo \`e mobile ($1$ byte);
  \item \textbf{posizione} --- coordinate correnti ($8$ byte);
  \item \textbf{timestamp} --- istante di trasmissione ($4$ byte);
  \item \textbf{lista dei vicini} --- $28$ byte per voce (id, timestamp e
  posizione); la distanza in hop di ogni voce \`e usata solo internamente e non
  pesa sulla dimensione on-air;
  \item \textbf{origin id} e \textbf{hop limit} --- $16+4$ byte, presenti solo in
  modalit\`a forwarded.
\end{itemize}
Il frame base \`e di $33$ byte e cresce di $28$ byte per ogni vicino annunciato.
La dimensione del beacon dipende quindi dalla densit\`a locale; questo aspetto
torna nel confronto fra le modalit\`a append e forwarded (Sezione~\ref{sec:multihop}).

\subsection{Operazione a livello MAC e topologia di rete}
\label{sec:macop}
L'accesso al canale segue la Distributed Coordination Function (DCF). Il nodo
rileva il canale; se \`e libero attende un DCF Interframe Space (DIFS,
$\SI{50}{\micro\second}$ nel modello), estrae un valore di backoff dalla
contention window $[0,CW-1]$ con $CW=16$ e $\text{slot}=\SI{20}{\micro\second}$, e
infine trasmette. Non si usano RTS/CTS n\'e ACK. Questo \`e adatto ai frame
broadcast e mantiene il protocollo semplice, ma il mittente non riceve alcun
riscontro sulla consegna, le collisioni non vengono rilevate e il problema del
terminale nascosto resta aperto. Da qui la necessit\`a di meccanismi adattivi.

Le boe formano un dominio ad hoc peer-to-peer (stile IBSS), senza access point,
router o server DHCP e senza associazioni persistenti. La rete \`e dinamica: i
nodi entrano ed escono dal raggio di comunicazione per correnti, onde e mobilit\`a.
Ogni nodo mantiene quindi una tabella dei vicini a stato leggero (\emph{soft
state}), aggiornata a ogni ricezione ed espirata periodicamente per togliere le
voci obsolete.

% ---------------------------------------------------------------------
\section{Static Beaconing Protocol (SBP)}
\label{sec:sbp}

L'SBP \`e il protocollo di riferimento. Ogni boa trasmette un beacon a intervallo
fisso $BI_{\text{static}}$, indipendentemente dalle condizioni di rete, ascolta in
continuazione, aggiorna la tabella dei vicini a ogni ricezione ed espira le voci
pi\`u vecchie di un timeout. \`E semplice e prevedibile. Il carico di canale e il
tasso di collisione, per\`o, crescono con il numero di nodi: in scenari densi le
collisioni riducono l'affidabilit\`a, mentre in scenari radi l'intervallo fisso
porta a trasmissioni inutilmente frequenti. Questi limiti motivano gli schemi
adattivi.

% ---------------------------------------------------------------------
\section{Adaptive Density-Aware Beaconing (ADAB)}
\label{sec:adab}

ADAB regola l'intervallo di beaconing in base a una sola grandezza locale: il
numero di vicini in tabella.

\subsection{Stima della densit\`a locale}
A ogni decisione di trasmissione il nodo conta i vicini correnti $N_t$ (dopo aver
espirato le voci obsolete) e li mappa su un fattore di densit\`a normalizzato
\begin{equation}
  f_t = \min\!\left(1,\; \frac{N_t}{\nthr}\right),
  \label{eq:adab-density}
\end{equation}
dove $\nthr$ \`e la soglia oltre la quale la congestione diventa rilevante.
Nell'implementazione $\nthr = 15$. Con $N_t = 0$ si ha $f_t = 0$; con
$N_t \ge \nthr$ si ha $f_t = 1$.

\subsection{Mapping quadratico dell'intervallo e jitter}
L'intervallo successivo $BI_{t+1}$ \`e una funzione quadratica del fattore di
densit\`a, interpolata fra un minimo e un massimo, con un piccolo jitter
moltiplicativo che desincronizza i nodi:
\begin{equation}
  BI_{t+1} = \Big[\,\bimin + f_t^{\,2}\,(\bimax-\bimin)\,\Big]\,(1+\epsilon_t),
  \qquad \epsilon_t \sim \mathcal{U}[-\eta,\eta],\;\; \eta=0.05,
  \label{eq:adab}
\end{equation}
con $BI_{t+1}$ vincolato all'intervallo $[\bimin,\bimax]$. Con pochi vicini
$f_t\approx 0$ e la boa trasmette vicino a $\bimin$, per una scoperta rapida; man
mano che $N_t$ si avvicina a $\nthr$ il termine quadratico porta l'intervallo
verso $\bimax$, riducendo le trasmissioni dove la contesa \`e maggiore. La forma
quadratica tiene il nodo reattivo a bassa densit\`a e ne accelera il rallentamento
ad alta densit\`a. I valori usati sono $\bimax = \SI{5.0}{\second}$ e $\bimin$
configurabile (negli esperimenti $\bimin \in \{0.25, 0.5, 1.0\}\,\si{\second}$).
ADAB cambia solo \emph{quando} i beacon vengono generati; il CSMA/CA sottostante
resta invariato.

% ---------------------------------------------------------------------
\section{Adaptive Context-Aware Beaconing (ACAB)}
\label{sec:acab}

\subsection{Motivazione}
ADAB reagisce alla sola densit\`a. La densit\`a, per\`o, non \`e l'unico segnale
locale utile a decidere se un beacon vada trasmesso. ACAB, lo scheduler proposto
in questa tesi, combina tre segnali: la densit\`a dei vicini, il tempo trascorso
dall'ultimo contatto e la mobilit\`a del nodo.

\subsection{Componente di densit\`a}
Come in ADAB, si calcola un conteggio normalizzato dei vicini
\begin{equation}
  s_{\mathrm{den}} = \min\!\left(1,\;\frac{N_t}{\nthr^{\mathrm{ACAB}}}\right),
  \qquad \nthr^{\mathrm{ACAB}} = 10 .
  \label{eq:acab-den}
\end{equation}

\subsection{Componente di contact-recency}
Se la boa ha ricevuto un beacon di recente, \`e meno urgente che si annunci; se \`e
rimasta a lungo in silenzio, conviene che trasmetta. Con
$\Delta = t - t_{\text{ultimo contatto}}$ e un orizzonte di recency
$\tau_c = \SI{20}{\second}$:
\begin{equation}
  s_{\mathrm{con}} =
  \begin{cases}
    \max\!\big(0,\; 1 - \Delta/\tau_c\big) & N_t>0 \text{ e contatto registrato},\\
    0 & \text{altrimenti.}
  \end{cases}
  \label{eq:acab-con}
\end{equation}
Un nodo contattato di recente d\`a $s_{\mathrm{con}}\to 1$ (rallenta); un nodo
isolato o silenzioso da tempo d\`a $0$ (trasmette pi\`u spesso).

\subsection{Componente di mobilit\`a}
Un nodo veloce cambia posizione rapidamente e dovrebbe aggiornarla pi\`u spesso.
Con velocit\`a $v=\lVert\mathbf{v}\rVert$ e una velocit\`a di riferimento
$v_{\text{ref}}$:
\begin{equation}
  s_{\mathrm{mob}} = \min\!\left(1,\;\frac{v}{v_{\text{ref}}}\right).
  \label{eq:acab-mob}
\end{equation}

\subsection{Combinazione pesata e calcolo dell'intervallo}
I tre punteggi sono combinati con pesi $(w_d,w_c,w_m)=(0.4,0.3,0.3)$ che sommano a
uno. La mobilit\`a riduce il backoff, da cui il termine $1-s_{\mathrm{mob}}$:
\begin{equation}
  s = w_d\,s_{\mathrm{den}} + w_c\,s_{\mathrm{con}} + w_m\,(1 - s_{\mathrm{mob}}).
  \label{eq:acab-blend}
\end{equation}
Il punteggio $s$ viene poi usato come $f_t$ nel mapping quadratico con jitter di
ADAB (Eq.~\ref{eq:adab}). In questo modo ACAB rallenta una boa densa, contattata
di recente e ferma, e tiene pi\`u frequente una boa isolata, silenziosa o veloce.
L'implementazione del calcolo \`e riportata nel Capitolo~\ref{ch:simulatore}.

\subsection{Confronto fra ADAB e ACAB}
ADAB e ACAB usano lo stesso mapping intervallo--punteggio e lo stesso jitter, e
differiscono solo nel modo in cui costruiscono il punteggio di backoff. ADAB usa
la sola densit\`a; ACAB ne fonde tre. Con ACAB un nodo denso ma in movimento non
viene frenato quanto un nodo altrettanto denso ma fermo, perch\'e la sua posizione
cambia; un nodo isolato che non sente nessuno da tempo trasmette in fretta a
prescindere dalla densit\`a. Il costo \`e un po' di stato e calcolo in pi\`u
(velocit\`a e ultimo contatto), trascurabile sui dispositivi considerati. Anche
ACAB agisce solo sul \emph{quando} trasmettere e lascia intatto il CSMA/CA.

% ---------------------------------------------------------------------
\section{Meccanismi di consapevolezza multihop}
\label{sec:multihop}

Un beacon trasporta consapevolezza a un solo hop. ProSafe la estende con due
meccanismi, selezionabili per esecuzione.

\subsection{Motivazione}
Un'imbarcazione che si avvicina a un gruppo di nuotatori ha interesse a conoscere
l'intero gruppo, compresi i membri momentaneamente nascosti da un'onda o da uno
scafo. La consapevolezza dovrebbe quindi propagarsi oltre il singolo hop radio;
ma ogni trasmissione spesa per propagarla compete per lo stesso canale, gi\`a
sotto pressione e privo di conferme. I due meccanismi seguenti sono due scelte
opposte nel compromesso fra copertura e costo di canale.

\subsection{Append mode (multihop passivo)}
In modalit\`a append una boa include, in ogni beacon, i nodi di cui \`e a
conoscenza: sia i vicini diretti (distanza $1$ hop) sia i nodi scoperti
indirettamente dalle liste dei vicini altrui. Alla ricezione, la boa integra la
lista annunciata dal mittente in una tabella di \emph{nodi scoperti},
incrementando di uno la distanza in hop di ogni voce. Non si invia nessun frame in
pi\`u: la consapevolezza viaggia sui beacon gi\`a previsti. Il costo \`e solo nella
dimensione del beacon ($28$ byte per nodo annunciato), e quindi nell'airtime per
frame.

Nell'implementazione attuale un eventuale limite di hop per l'append
(\texttt{append\_hop\_limit}) non \`e applicato: una boa annuncia tutti i nodi
scoperti, a prescindere dalla distanza in hop. La distanza in hop resta comunque
registrata e potrebbe servire in futuro a limitare la propagazione.

\subsection{Forwarded mode (multihop attivo)}
In modalit\`a forwarded una boa ritrasmette i beacon che riceve. Ogni beacon porta
un \emph{origin id} (il mittente originale) e un \emph{hop limit} (un contatore
TTL). Quando riceve un beacon che non ha originato e con TTL positivo, la boa pu\`o
accodare una copia da rilanciare con il TTL decrementato. Origine, timestamp e
payload restano invariati, cos\`i da poter riconoscere i duplicati ed evitare che
un nodo rilanci un beacon da s\'e originato. I rilanci passano per la stessa
pipeline CSMA dei beacon propri, dopo un piccolo jitter casuale che desincronizza
i molti ricevitori dello stesso frame.

\subsection{Deduplicazione, coda di forwarding e prevenzione dei loop}
Senza controlli, il forwarding attivo pu\`o generare un \emph{broadcast storm}.
L'implementazione adotta tre meccanismi di contenimento:
\begin{itemize}[leftmargin=1.4em,nosep]
  \item \textbf{Deduplicazione per origine}: ogni boa memorizza, per ciascun
  \emph{origin id}, il timestamp pi\`u recente gi\`a inoltrato; un beacon viene
  accodato solo se pi\`u fresco dell'ultimo inoltrato dalla stessa origine. Questo
  evita di rilanciare pi\`u volte la stessa informazione e spezza i cicli.
  \item \textbf{Limite del TTL}: il campo \emph{hop limit}, decrementato a ogni
  rilancio, limita la profondit\`a di propagazione (valore tipico $1$ hop).
  \item \textbf{Coda limitata}: i rilanci pendenti stanno in una coda di
  dimensione massima \texttt{pending\_queue\_limit}; quando \`e piena i nuovi
  beacon vengono scartati senza essere registrati, cos\`i che una copia successiva
  possa entrare se nel frattempo si libera spazio.
\end{itemize}
In sintesi, la modalit\`a forwarded corrente accoda e rilancia ogni beacon fresco,
entro TTL e non proprio, con il solo limite della coda.

\subsection{Copertura e costo di canale}
I due meccanismi sono scelte opposte. L'append \`e economico (nessun frame in
pi\`u) ma limitato dalla freschezza e dalla dimensione crescente del beacon; il
forwarded estende la portata fisica ma aumenta traffico, collisioni e latenza.
Quale dei due convenga in una rete broadcast densa, a dominio unico e senza
conferme, \`e valutato nel Capitolo~\ref{ch:risultati}.

% ---------------------------------------------------------------------
\section*{Sintesi del capitolo}
Questo capitolo ha presentato l'architettura ProSafe e i protocolli di beaconing.
L'SBP \`e una linea di base semplice ma non scalabile; ADAB adatta l'intervallo
alla densit\`a locale; ACAB estende l'adattamento combinando densit\`a, recency
del contatto e mobilit\`a. Sono stati inoltre descritti i due meccanismi di
consapevolezza multihop --- append (passivo) e forwarded (attivo), con i relativi
meccanismi di deduplicazione e contenimento. Il prossimo capitolo descrive il
simulatore sviluppato per valutare questi protocolli in condizioni controllate e
riproducibili.

% =====================================================================
%  BOZZA - Capitolo 4 (italiano)
%  Risultati: protocolli di beaconing e valutazione sperimentale.
%  Le sezioni su SBP, ADAB e ACAB provengono dal Capitolo 2 e sono
%  state riviste nel registro; dalla fine della Sezione "Protocolli di
%  beaconing" in poi prosegue l'autore con la valutazione.
% =====================================================================

\chapter{Risultati}
\label{ch:risultati}

Questo capitolo presenta la valutazione sperimentale delle modalit\`a multi-hop
introdotte nel Capitolo~\ref{ch:contributo}. Poich\'e ogni configurazione
di simulazione combina una delle due modalit\`a con uno dei protocolli di
beaconing che regolano gli intervalli di trasmissione, il capitolo si apre con
la descrizione dei tre protocolli --- SBP, ADAB e ACAB --- gi\`a esistenti nel simulatore.
Seguono la descrizione dello scenario di rete adottato negli esperimenti
(Sezione~\ref{sec:scenario}), la definizione delle metriche impiegate nella
valutazione (Sezione~\ref{sec:eval-metrics}) e l'analisi delle prestazioni
ottenute nelle simulazioni.

% ---------------------------------------------------------------------
\section{Protocolli di beaconing}
\label{sec:beaconing}
I tre protocolli agiscono esclusivamente sull'istante di generazione del
beacon: una volta generato, il frame attraversa la pipeline di trasmissione
CSMA/CA descritta nel Capitolo~\ref{ch:simulatore}. Ciascun protocollo pu\`o
essere combinato con una qualunque delle modalit\`a multi-hop del
Capitolo~\ref{ch:contributo}.

\subsection{Static Beaconing Protocol (SBP)}
\label{sec:sbp}
Lo SBP costituisce il protocollo di riferimento: ogni boa trasmette un beacon a
intervallo fisso $BI_{\text{static}}$, indipendentemente dalle condizioni della
rete. La semplicit\`a e la prevedibilit\`a del comportamento ne fanno la base di
confronto naturale; tuttavia, negli scenari densi l'assenza di adattamento
accresce la probabilit\`a di collisione e degrada l'affidabilit\`a, mentre negli
scenari sparsi l'intervallo fisso produce trasmissioni pi\`u frequenti del
necessario. Tali limiti motivano l'adozione di protocolli adattivi.

\subsection{Adaptive Density-Aware Beaconing (ADAB)}
\label{sec:adab}
ADAB regola l'intervallo sulla base di un'unica grandezza locale: il numero di
vicini presenti in tabella. A ogni decisione di scheduling il nodo rileva il
numero di vicini correnti $N_t$ e ne deriva un fattore di densit\`a normalizzato
\begin{equation}
  f_t = \min\!\left(1,\; \frac{N_t}{\nthr}\right),
  \qquad \nthr = 15,
  \label{eq:adab-density}
\end{equation}
dove $\nthr$ \`e la soglia oltre la quale la congestione \`e considerata
rilevante. L'intervallo successivo \`e ottenuto da una funzione quadratica del
fattore di densit\`a, interpolata fra un minimo e un massimo, con un jitter
moltiplicativo $\delta$ che desincronizza le trasmissioni dei nodi:
\begin{equation}
  BI_{t+1} = \Big[\,\bimin + f_t^{\,2}\,(\bimax-\bimin)\,\Big]\,(1+\delta),
  \qquad \delta \sim \mathcal{U}[-\eta,\eta],\;\; \eta=0.05,
  \label{eq:adab}
\end{equation}
con $BI_{t+1}$ vincolato all'intervallo $[\bimin,\bimax]$. Con pochi vicini
$f_t\approx 0$ e la boa trasmette a intervalli prossimi a $\bimin$, favorendo
una scoperta rapida; man mano che $N_t$ si avvicina a $\nthr$, il termine
quadratico spinge l'intervallo verso $\bimax$, riducendo le trasmissioni dove la
contesa \`e maggiore. La dipendenza quadratica mantiene quindi il nodo reattivo
alle basse densit\`a e ne accentua il rallentamento all'aumentare della
congestione. I valori adottati sono $\bimax = \SI{5.0}{\second}$ e $\bimin$
configurabile (negli esperimenti $\bimin \in \{0.25, 0.5, 1.0\}\,\si{\second}$).

\subsection{Adaptive Context-Aware Beaconing (ACAB)}
\label{sec:acab}
La densit\`a non \`e l'unico segnale locale utile a stabilire l'opportunit\`a di
una trasmissione. ACAB combina tre segnali: la densit\`a dei vicini, il tempo
trascorso dall'ultimo contatto e la mobilit\`a del nodo.

\paragraph{Componente di densit\`a.} Come in ADAB, si calcola un conteggio
normalizzato dei vicini, con una soglia pi\`u bassa:
\begin{equation}
  s_{\mathrm{den}} = \min\!\left(1,\;\frac{N_t}{\nthr^{\mathrm{ACAB}}}\right),
  \qquad \nthr^{\mathrm{ACAB}} = 10 .
  \label{eq:acab-den}
\end{equation}

\paragraph{Componente di tempo dall'ultimo contatto.} La ricezione recente di un
beacon riduce l'urgenza di trasmettere; un periodo prolungato privo di contatti
la accresce. Detto $\Delta = t - t_{\text{ultimo contatto}}$ il tempo trascorso
dall'ultimo beacon ricevuto e $\tau_c = \SI{20}{\second}$ la durata oltre la
quale un contatto non \`e pi\`u considerato recente:
\begin{equation}
  s_{\mathrm{con}} =
  \begin{cases}
    \max\!\big(0,\; 1 - \Delta/\tau_c\big) & N_t>0 \text{ e contatto registrato},\\
    0 & \text{altrimenti.}
  \end{cases}
  \label{eq:acab-con}
\end{equation}
Un nodo contattato di recente produce $s_{\mathrm{con}}\to 1$ (e rallenta); un
nodo isolato o a lungo silenzioso produce $0$ (e trasmette pi\`u di frequente).

\paragraph{Componente di mobilit\`a.} Un nodo in rapido movimento invalida in
breve tempo l'informazione di posizione annunciata e deve pertanto aggiornarla
con maggiore frequenza. Con velocit\`a $v=\lVert\mathbf{v}\rVert$ e velocit\`a
di riferimento $v_{\text{ref}}$ (la velocit\`a di default, pari a
\SI{15}{\metre\per\second}):
\begin{equation}
  s_{\mathrm{mob}} = \min\!\left(1,\;\frac{v}{v_{\text{ref}}}\right).
  \label{eq:acab-mob}
\end{equation}

\paragraph{Combinazione e calcolo dell'intervallo.} I tre punteggi sono
combinati mediante pesi definiti nella configurazione,
$(w_d,w_c,w_m)=(0.4,0.3,0.3)$, che ne determinano il contributo relativo. La
mobilit\`a riduce il backoff, da cui il termine $1-s_{\mathrm{mob}}$:
\begin{equation}
  s = w_d\,s_{\mathrm{den}} + w_c\,s_{\mathrm{con}} + w_m\,(1 - s_{\mathrm{mob}}).
  \label{eq:acab-blend}
\end{equation}
Il punteggio $s$ sostituisce $f_t$ nell'espressione dell'intervallo di ADAB
(Eq.~\ref{eq:adab}): ACAB rallenta cos\`i una boa densa, contattata di recente e
ferma, e mantiene pi\`u frequente una boa isolata, silenziosa o in movimento.
Rispetto ad ADAB, che impiega la sola densit\`a, ACAB costruisce il punteggio di
backoff a partire da tre segnali e richiede quindi il mantenimento di un numero
maggiore di variabili di stato; l'overhead computazionale resta comunque
trascurabile per i dispositivi considerati.

% ---------------------------------------------------------------------
\section{Dinamicit\`a della rete}
\label{sec:scenario}
La dinamicit\`a della rete dipende dallo \texttt{scenario}
configurato (\emph{static}, \emph{ramp}, \emph{random}), che determina l'evoluzione del
numero di nodi nel corso della simulazione --- meccanismo descritto nella Sezione~\ref{sec:dynamics} del Capitolo~\ref{ch:simulatore}. 
Ai fini della valutazione sperimentale \`e
stato adottato lo scenario random: la popolazione di nodi cresce e decresce periodicamente
di una frazione percentuale definita dalla configurazione, cos\`i da riprodurre il
comportamento di una rete reale, in cui i nodi entrano ed escono di continuo dal raggio di
copertura radio.

% ---------------------------------------------------------------------
% Da qui in poi: valutazione delle prestazioni delle simulazioni
% (sezioni a cura dell'autore).
% ---------------------------------------------------------------------

% ---------------------------------------------------------------------
\section{Metriche di valutazione}
\label{sec:eval-metrics}
Questa sezione definisce le metriche con cui sono valutate le prestazioni
delle configurazioni sperimentali. I valori sono calcolati al termine di ogni
simulazione a partire dai dati grezzi accumulati dalla classe \texttt{Metrics}
(Sezione~\ref{sec:metrics}); per ciascuna metrica si indica dapprima che cosa
esprime e quindi la formula del calcolo effettuato.

\paragraph{Percentuale media di rete scoperta.} Esprime la conoscenza media
della topologia di rete, ovvero quanta parte della rete conosce in media ciascuna boa. 
Ogni boa $i$ accumula l'insieme
$S_i$ dei nodi \emph{unici} di cui \`e venuta a conoscenza:  con ogni beacon
ricevuto si uniscono ad esso il mittente, l'eventuale origine (in modalit\`a
forwarding singolo) e tutti i nodi della lista annunciata. Detta $N$ la densit\`a della rete:
\begin{equation}
  \overline{U} = \frac{1}{N} \sum_{i=1}^{N} |S_i|,
  \qquad
  D_{\%} = \frac{\overline{U}}{N-1} \cdot 100,
  \label{eq:netdisc}
\end{equation}
dove il denominatore $N-1$ \`e il massimo numero di nodi che una boa pu\`o
scoprire, poich\'e $S_i$ non include mai la boa stessa. A scoperta completa
ogni $|S_i|$ vale $N-1$, dunque $\overline{U}=N-1$ e $D_{\%}=100\%$.

\paragraph{B-PDR (Packet Delivery Ratio per broadcast).} Esprime
l'affidabilit\`a delle trasmissioni, ossia la probabilit\`a che un frame
in aria raggiunga integro un nodo in ascolto. Per ogni trasmissione, il canale
conta i destinatari entro la portata (opportunit\`a di ricezione) e le boe
contano le ricezioni andate a buon fine:
\begin{equation}
  \text{B-PDR} = \frac{R_{\text{rx}}}{R_{\text{opp}}},
  \label{eq:bpdr}
\end{equation}
con $R_{\text{rx}}$ il totale delle ricezioni riuscite e $R_{\text{opp}}$ il
totale delle opportunit\`a (trasmissione $\times$ destinatario in portata).

\paragraph{Collision rate.} Esprime la frazione delle opportunit\`a di
ricezione $R_{\text{opp}}$ perse per collisione:
\begin{equation}
  CR = \frac{C}{R_{\text{opp}}},
  \label{eq:collrate}
\end{equation}
con $C$ il totale delle collisioni avvenute, sia dirette che provocate da un terminale nascosto (Sezione~\ref{sec:broadcast-impl}).

\paragraph{Packet loss.} Esprime la frazione complessiva delle opportunit\`a
di ricezione $R_{\text{opp}}$ perse a causa di collisioni ed errori del canale non ideale (Eq.~\ref{eq:delivery-prob}):
\begin{equation}
  PL = \frac{C + E_{\text{canale}}}{R_{\text{opp}}},
  \label{eq:packetloss}
\end{equation}
con $E_{\text{canale}}$ le perdite dovute al canale. Vale $PL \ge CR$, e la
differenza tra le due evidenzia il contributo delle perdite dovute al canale
rispetto a quelle della contesa.

\paragraph{Latenze.} Esprimono i tempi relativi alla diffusione dei
beacon nella rete. Vengono registrate la \emph{latenza di consegna} (tempo tra
generazione e prima ricezione di un beacon), la \emph{latenza di reazione}
(tempo perch\'e un nodo scopra per la prima volta una data origine) e la
gi\`a citata latenza di scheduler (tempo trascorso tra la decisione di trasmettere e l’effettivo inizio della trasmissione, Sezione~\ref{sec:txpipeline}).

% ---------------------------------------------------------------------
\section{Piano degli esperimenti}
\label{sec:eval-setup}
La valutazione confronta le tre modalit\`a di inoltro --- \emph{single-hop}
(nessun inoltro, base di confronto), \emph{forwarding aggregato} (append) e
\emph{forwarding singolo} (forward) --- al variare della densit\`a della rete e
delle condizioni del canale. Ogni configurazione sperimentale \`e identificata
dalla quadrupla (modalit\`a multihop, protocollo di beaconing, intervallo
minimo di beaconing, condizione del canale); per ciascuna quadrupla la
densit\`a varia da 20 a 100 boe a passi di 20, in un'area di
$500 \times 500\,\si{\metre\squared}$, con scenario \emph{random}
(variabilit\`a del 10\% per aggiornamento) e durata simulata di
\SI{120}{\second}. Ogni punto \`e la media di 10 repliche indipendenti con
semi diversi; dove indicato, le bande ombreggiate nei grafici temporali
rappresentano la deviazione standard tra le repliche. La campagna qui
analizzata adotta SBP come protocollo di beaconing: il suo intervallo fisso
rende direttamente confrontabile il costo delle modalit\`a di inoltro, senza
che un adattamento dell'intervallo mascheri gli effetti della contesa.
% NOTA: se verranno completate le campagne con ADAB/ACAB, estendere qui
% e reintrodurre i protocolli adattivi nella griglia sperimentale.

Le condizioni del canale sono due. Nel \emph{canale ideale}
(\texttt{ideal\_channel = true}) ogni frame che raggiunge un nodo in portata
viene consegnato, salvo collisione: le uniche perdite derivano dalla contesa,
quindi $PL = CR$. Nel \emph{canale non ideale}
(\texttt{ideal\_channel = false}) la consegna \`e probabilistica secondo il
modello a due soglie dell'Eq.~\ref{eq:delivery-prob} (probabilit\`a 0.9 entro
\SI{70}{\metre}, 0.15 fino a \SI{120}{\metre}): alle perdite da collisione si
sommano gli errori di canale e la differenza $PL - CR$ ne misura il
contributo. Gli intervalli minimi di beaconing considerati sono
$BI \in \{1.0, 0.5, 0.25\}\,\si{\second}$. La
Tabella~\ref{tab:eval-grid} riassume la griglia sperimentale.

\begin{table}[htbp]
  \centering
  \caption[Griglia dei parametri sperimentali]{Griglia dei parametri della
  campagna sperimentale. Ogni combinazione \`e mediata su 10 repliche.}
  \label{tab:eval-grid}
  \begin{tabular}{ll}
    \toprule
    \textbf{Parametro} & \textbf{Valori} \\
    \midrule
    Modalit\`a multihop & single-hop, append, forward (limite 4 hop) \\
    Protocollo di beaconing & SBP \\
    Intervallo di beaconing $BI$ & 1.0, 0.5, \SI{0.25}{\second} \\
    Densit\`a (numero di boe) & 20, 40, 60, 80, 100 \\
    Canale & ideale, non ideale \\
    Scenario & random (variabilit\`a 10\%) \\
    Durata / repliche & \SI{120}{\second} / 10 semi \\
    \bottomrule
  \end{tabular}
\end{table}

\paragraph{Lettura dei grafici di confronto.} Nei grafici di confronto tra
modalit\`a, le barre raggruppate per densit\`a confrontano single-hop, append
e forward sulla stessa metrica; i grafici temporali mostrano invece
l'evoluzione della scoperta della rete nella configurazione pi\`u densa
(100 boe). I
risultati sono presentati prima per il canale ideale
(Sezione~\ref{sec:res-ideal}), che isola l'effetto della contesa, poi per il
canale non ideale (Sezione~\ref{sec:res-error}), che aggiunge le perdite
probabilistiche; seguono l'analisi dell'effetto dell'intervallo di beaconing
(Sezione~\ref{sec:res-interval}) e la discussione complessiva
(Sezione~\ref{sec:discussione}).

% ---------------------------------------------------------------------
\section{Risultati con canale ideale}
\label{sec:res-ideal}
In assenza di errori di canale ogni perdita \`e imputabile alla contesa:
questa condizione isola il costo intrinseco delle modalit\`a di inoltro, che
accrescono il traffico (forward) o la dimensione dei frame (append) rispetto
al beaconing single-hop.

\subsection{Scoperta della rete}
\label{sec:res-ideal-disc}
% Figure suggerite (cartella "senza errore"):
%   mode_comparison_avg_percentage_network_discovered_interval-<tag>.png
%   mode_comparison_discovery_growth_interval-<tag>.png
\begin{figure}[htbp]
  \centering
  % TODO: \includegraphics[width=\textwidth]{main/images/ideal/mode_comparison_avg_percentage_network_discovered_interval-1_0.png}
  \caption[Percentuale di rete scoperta, canale ideale]{Percentuale media di
  rete scoperta al variare della densit\`a per le tre modalit\`a di inoltro
  (canale ideale, $BI=\SI{1.0}{\second}$).}
  \label{fig:ideal-discovery}
\end{figure}
La Figura~\ref{fig:ideal-discovery} mostra la percentuale media di rete
scoperta ($D_{\%}$, Eq.~\ref{eq:netdisc}). Il single-hop \`e limitato dalla
portata radio: una boa scopre soltanto i nodi da cui riceve direttamente,
perci\'o $D_{\%}$ cresce con la densit\`a solo per effetto della maggiore
connettivit\`a locale, passando da circa l'81\% a 20 boe a circa il 91\% a
100 boe senza mai raggiungere la copertura completa. Le due modalit\`a
multi-hop estendono la conoscenza oltre la portata: append veicola
l'informazione dei vicini all'interno dei beacon periodici, forward la
propaga con ritrasmissioni esplicite. Entrambe saturano la metrica:
gi\`a a 20 boe superano il 99\% e da 40 boe in su raggiungono di fatto il
100\% a ogni densit\`a. Il divario rispetto al single-hop \`e massimo nella
configurazione pi\`u sparsa ($\approx$19 punti percentuali a 20 boe), dove i
collegamenti diretti sono pochi e l'informazione di secondo e terzo hop \`e
determinante, e si riduce a $\approx$9 punti alle densit\`a maggiori.

La dinamica temporale della scoperta \`e riportata nella
Figura~\ref{fig:ideal-growth} per la configurazione pi\`u densa.
\begin{figure}[htbp]
  \centering
  % TODO: \includegraphics[width=\textwidth]{main/images/ideal/mode_comparison_discovery_growth_interval-1_0.png}
  \caption[Crescita della scoperta nel tempo, canale ideale]{Evoluzione
  temporale della percentuale di rete scoperta nella configurazione a 100 boe
  (canale ideale, $BI=\SI{1.0}{\second}$); le bande indicano la deviazione
  standard tra le repliche.}
  \label{fig:ideal-growth}
\end{figure}
Le tre curve differiscono soprattutto nella velocit\`a di convergenza.
Append raggiunge il 100\% in circa 3--4 secondi: sono sufficienti pochi cicli
di beaconing perch\'e le liste di vicinato annunciate propaghino la
conoscenza dell'intera topologia. Forward sale anch'esso rapidamente
($\approx$97\% dopo 5 secondi) ma completa la scoperta pi\`u gradualmente,
intorno ai 20 secondi: ogni ritrasmissione veicola un solo beacon di
origine, quindi la copertura degli ultimi nodi richiede pi\`u scambi. Il
single-hop, privo di meccanismi di estensione, converge lentamente verso un
plateau parziale: supera l'80\% solo dopo circa 30 secondi e si assesta
attorno al 91--92\% a fine simulazione. Le bande di deviazione standard sono
pressoch\'e nulle per le modalit\`a multi-hop e visibilmente pi\`u ampie per
il single-hop, la cui scoperta dipende dagli incontri casuali dettati dalla
mobilit\`a delle singole repliche.

\subsection{Affidabilit\`a e contesa}
\label{sec:res-ideal-pdr}
% Figure suggerite (cartella "senza errore"):
%   mode_comparison_pdr_interval-<tag>.png
%   mode_comparison_collision_rate_interval-<tag>.png
\begin{figure}[htbp]
  \centering
  % TODO: \includegraphics[width=\textwidth]{main/images/ideal/mode_comparison_pdr_interval-1_0.png}
  \caption[B-PDR, canale ideale]{B-PDR al variare della densit\`a per le tre
  modalit\`a di inoltro (canale ideale, $BI=\SI{1.0}{\second}$).}
  \label{fig:ideal-pdr}
\end{figure}
\begin{figure}[htbp]
  \centering
  % TODO: \includegraphics[width=\textwidth]{main/images/ideal/mode_comparison_collision_rate_interval-1_0.png}
  \caption[Collision rate, canale ideale]{Collision rate al variare della
  densit\`a (canale ideale, $BI=\SI{1.0}{\second}$).}
  \label{fig:ideal-collision}
\end{figure}
La maggiore conoscenza della rete ha un costo in termini di contesa. In
modalit\`a forward ogni beacon ricevuto pu\`o generare una ritrasmissione:
il traffico offerto cresce e con esso la probabilit\`a di collisione,
specialmente alle densit\`a elevate. In modalit\`a append il numero di
trasmissioni resta quello del beaconing periodico, ma i frame pi\`u lunghi
occupano il canale pi\`u a lungo, accrescendo la finestra di vulnerabilit\`a
alle collisioni.

I dati confermano questa gerarchia di costo. Il single-hop mantiene un B-PDR
superiore a 0.93 a ogni densit\`a, con un collision rate che non supera 0.07.
Append degrada in modo graduale: da 0.99 a 20 boe a circa 0.74 a 100 boe,
con collision rate 0.27 nel caso peggiore. Forward crolla invece gi\`a alle
densit\`a intermedie: 0.68 a 40 boe, 0.44 a 60 e appena 0.26 a 100 boe, dove
quasi tre opportunit\`a di ricezione su quattro ($CR \approx 0.74$) sono
perse per collisione. Le ritrasmissioni moltiplicano il traffico offerto in
proporzione alla densit\`a, innescando un circolo vizioso: pi\`u vicini
significano pi\`u inoltri, quindi pi\`u collisioni. Nel canale ideale packet
loss e collision rate coincidono ($PL = CR$, come atteso dall'assenza di
errori di canale), e infatti in ogni configurazione risulta
$\text{B-PDR} + PL \approx 1$: il confronto tra le
Figure~\ref{fig:ideal-pdr} e~\ref{fig:ideal-collision} lo verifica come
controllo di coerenza dei dati. Append rappresenta quindi il miglior
compromesso tra conoscenza e affidabilit\`a: ottiene la stessa scoperta di
forward a una frazione del suo costo in collisioni.

\subsection{Efficienza informativa dei beacon}
\label{sec:res-ideal-ratio}
% Figura suggerita (cartella "senza errore"):
%   mode_comparison_neighbor_receiver_ratio_interval-<tag>.png
\begin{figure}[htbp]
  \centering
  % TODO: \includegraphics[width=\textwidth]{main/images/ideal/mode_comparison_neighbor_receiver_ratio_interval-1_0.png}
  \caption[Rapporto vicini/riceventi, canale ideale]{Rapporto tra nodi
  annunciati nel beacon e riceventi in portata al momento della trasmissione
  (canale ideale, $BI=\SI{1.0}{\second}$); la linea tratteggiata orizzontale
  indica la parit\`a (rapporto = 1), le linee tratteggiate per modalit\`a il
  delta medio vicini--riceventi.}
  \label{fig:ideal-ratio}
\end{figure}
Il rapporto tra i nodi annunciati in un beacon e i riceventi effettivamente in
portata quantifica l'\emph{amplificazione informativa} dell'inoltro: valori
superiori a 1 indicano che il beacon trasporta conoscenza relativa a nodi che
il mittente non pu\`o raggiungere direttamente, ovvero informazione acquisita
per via multi-hop.

In modalit\`a append il rapporto cresce da circa 4.5 a 20 boe a 6 a 100 boe:
ogni beacon annuncia in media sei volte pi\`u nodi di quanti possano
riceverlo, e il delta medio vicini--riceventi passa da $+10$ a $+80$,
avvicinandosi al limite teorico $N-1$ meno i riceventi in portata --- segno
che le liste annunciate arrivano a coprire l'intera rete. Single-hop e
forward restano invece prossimi alla parit\`a (rapporto $\approx$1.3, delta
compreso tra $+1$ e $+5$): il primo perch\'e annuncia solo i vicini diretti,
il secondo perch\'e l'informazione aggiuntiva viaggia in beacon ritrasmessi
separatamente anzich\'e in liste aggregate. L'amplificazione informativa di
append \`e dunque ottenuta a parit\`a di numero di trasmissioni, al solo
costo di frame pi\`u lunghi.

% ---------------------------------------------------------------------
\section{Risultati con canale non ideale}
\label{sec:res-error}
Il canale non ideale sovrappone alle collisioni le perdite probabilistiche
del modello a due soglie: i beacon oltre \SI{70}{\metre} vengono consegnati
con probabilit\`a 0.15, il che riduce drasticamente l'affidabilit\`a dei
collegamenti lunghi --- proprio quelli su cui l'inoltro multi-hop costruisce
il proprio vantaggio. Questa sezione ripercorre le metriche della
Sezione~\ref{sec:res-ideal} nelle stesse configurazioni, con l'obiettivo di
valutare la robustezza delle modalit\`a di inoltro in condizioni realistiche.

\subsection{Scoperta della rete}
\label{sec:res-error-disc}
\begin{figure}[htbp]
  \centering
  % TODO: \includegraphics[width=\textwidth]{main/images/error/mode_comparison_avg_percentage_network_discovered_interval-1_0.png}
  \caption[Percentuale di rete scoperta, canale non ideale]{Percentuale media
  di rete scoperta al variare della densit\`a (canale non ideale,
  $BI=\SI{1.0}{\second}$).}
  \label{fig:error-discovery}
\end{figure}
Il confronto con la Figura~\ref{fig:ideal-discovery} mostra che il degrado si
concentra sul single-hop, la cui scoperta cala di 5--7 punti percentuali a
ogni densit\`a (dal 76\% a 20 boe all'84--85\% a 100 boe): ogni beacon perso
dal canale \`e per esso un'occasione di scoperta irrimediabilmente persa, e
la copertura resta lontana dal valore gi\`a parziale del caso ideale. Le
modalit\`a multi-hop assorbono invece quasi completamente le perdite: a
parte la configurazione pi\`u sparsa (20 boe, dove append si ferma attorno
al 94\%), da 40 boe in su entrambe restano sopra il 99\%. La ridondanza
dell'inoltro --- la stessa informazione raggiunge un nodo lungo percorsi e in
istanti diversi --- compensa la perdita probabilistica dei singoli frame, e
il vantaggio sul single-hop si amplifica rispetto al canale ideale
(fino a $\approx$18 punti a 20 boe e $\approx$15 punti a 100 boe).
\begin{figure}[htbp]
  \centering
  % TODO: \includegraphics[width=\textwidth]{main/images/error/mode_comparison_discovery_growth_interval-1_0.png}
  \caption[Crescita della scoperta nel tempo, canale non ideale]{Evoluzione
  temporale della percentuale di rete scoperta nella configurazione a 100 boe
  (canale non ideale, $BI=\SI{1.0}{\second}$).}
  \label{fig:error-growth}
\end{figure}
La dinamica temporale (Figura~\ref{fig:error-growth}) conserva l'ordinamento
del caso ideale ma ne dilata i tempi: append completa la scoperta in circa
5--6 secondi (contro 3--4), forward raggiunge il 90\% intorno ai 10 secondi
e si avvicina asintoticamente al 100\% solo verso fine simulazione, mentre il
single-hop cresce lentamente fino all'84--85\% a 120 secondi, senza plateau
evidente. Le perdite del canale allungano quindi la coda della scoperta ---
servono pi\`u tentativi perch\'e un beacon superi i collegamenti a bassa
probabilit\`a --- ma non ne cambiano l'esito per le modalit\`a multi-hop.

\subsection{Affidabilit\`a, collisioni e perdite}
\label{sec:res-error-pdr}
\begin{figure}[htbp]
  \centering
  % TODO: \includegraphics[width=\textwidth]{main/images/error/mode_comparison_pdr_interval-1_0.png}
  \caption[B-PDR, canale non ideale]{B-PDR al variare della densit\`a
  (canale non ideale, $BI=\SI{1.0}{\second}$).}
  \label{fig:error-pdr}
\end{figure}
\begin{figure}[htbp]
  \centering
  % TODO: \includegraphics[width=\textwidth]{main/images/error/mode_comparison_loss_rate_interval-1_0.png}
  \caption[Packet loss, canale non ideale]{Packet loss complessivo al variare
  della densit\`a (canale non ideale, $BI=\SI{1.0}{\second}$).}
  \label{fig:error-loss}
\end{figure}
Nel canale non ideale il packet loss (Eq.~\ref{eq:packetloss}) scompone le
perdite in due contributi: la contesa ($CR$) e gli errori di canale
($PL - CR$). Il confronto tra la Figura~\ref{fig:error-loss} e il collision
rate consente di attribuire il degrado a ciascuna causa.

Il modello di canale impone anzitutto un tetto all'affidabilit\`a,
indipendente dalla modalit\`a: gi\`a a 20 boe, dove la contesa \`e
trascurabile, il B-PDR di tutte le configurazioni si ferma a 0.42--0.45 e il
packet loss parte da una base di circa 0.55--0.58. Tale valore riflette la
geometria delle opportunit\`a di ricezione: la maggior parte dei riceventi in
portata cade nella corona esterna (70--\SI{120}{\metre}), dove la
probabilit\`a di consegna \`e appena 0.15. Su questa base i comportamenti
divergono con la densit\`a. Il single-hop resta pressoch\'e piatto (B-PDR
$\approx$0.41 e $PL \approx 0.59$ a 100 boe): le sue perdite sono dominate
dal canale, non dalla contesa ($CR \le 0.03$). Append aggiunge una contesa
moderata ($CR \approx 0.18$ a 100 boe, B-PDR 0.34), mentre forward somma al
tetto del canale il proprio traffico di ritrasmissione: $CR$ sale a 0.56 e il
packet loss complessivo raggiunge 0.81, con B-PDR 0.19. \`E interessante
notare che il collision rate di forward risulta inferiore a quello del caso
ideale a parit\`a di densit\`a (0.56 contro 0.74): i beacon persi dal canale
non vengono ricevuti e dunque nemmeno rilanciati, sicch\'e l'errore di canale
agisce indirettamente da freno sul traffico di inoltro.

\subsection{Efficienza informativa dei beacon}
\label{sec:res-error-ratio}
\begin{figure}[htbp]
  \centering
  % TODO: \includegraphics[width=\textwidth]{main/images/error/mode_comparison_neighbor_receiver_ratio_interval-1_0.png}
  \caption[Rapporto vicini/riceventi, canale non ideale]{Rapporto tra nodi
  annunciati e riceventi in portata (canale non ideale,
  $BI=\SI{1.0}{\second}$).}
  \label{fig:error-ratio}
\end{figure}
Il quadro \`e analogo al canale ideale, con valori lievemente ridotti:
append cresce da un rapporto di circa 2.6 a 20 boe a 5.9 a 100 boe (delta da
$+4.8$ a $+77$), mentre single-hop e forward restano attorno alla parit\`a.
La riduzione rispetto al caso ideale (a 20 boe il rapporto di append passa da
4.5 a 2.6) riflette il pi\`u lento popolamento delle tabelle dei vicini: le
perdite del canale ritardano sia l'acquisizione dell'informazione da
annunciare sia, per il single-hop, il mantenimento stesso della tabella ---
il cui delta diventa lievemente negativo ($-0.7$ a 100 boe), segno di
tabelle che sottostimano i riceventi effettivamente in portata.

% ---------------------------------------------------------------------
\section{Effetto dell'intervallo di beaconing}
\label{sec:res-interval}
Le sezioni precedenti hanno fissato $BI=\SI{1.0}{\second}$; questa sezione
analizza la sensibilit\`a dei risultati all'intervallo, ripetendo i confronti
con $BI=\SI{0.5}{\second}$ e $BI=\SI{0.25}{\second}$. Intervalli pi\`u brevi
accelerano la scoperta e riducono la latenza di reazione, ma moltiplicano il
traffico offerto e quindi la contesa. La Tabella~\ref{tab:interval-summary}
riassume scoperta e affidabilit\`a nella configurazione pi\`u densa.

\begin{table}[htbp]
  \centering
  \caption[Sensibilit\`a all'intervallo di beaconing (100 boe)]{Percentuale
  di rete scoperta e B-PDR a 100 boe al variare dell'intervallo di beaconing,
  nei due regimi di canale. Valori medi su 10 repliche.}
  \label{tab:interval-summary}
  \begin{tabular}{llcccccc}
    \toprule
    & & \multicolumn{2}{c}{\textbf{Single-hop}} & \multicolumn{2}{c}{\textbf{Append}} & \multicolumn{2}{c}{\textbf{Forward}} \\
    \cmidrule(lr){3-4}\cmidrule(lr){5-6}\cmidrule(lr){7-8}
    \textbf{Canale} & $BI$ (\si{\second}) & $D_{\%}$ & B-PDR & $D_{\%}$ & B-PDR & $D_{\%}$ & B-PDR \\
    \midrule
    ideale     & 1.0  & $\sim$91 & 0.93 & 100 & 0.74 & 100      & 0.26 \\
    ideale     & 0.5  & $\sim$91 & 0.87 & 100 & 0.51 & $\sim$99.5 & 0.21 \\
    ideale     & 0.25 & $\sim$91 & 0.75 & 100 & 0.23 & $\sim$98.5 & 0.18 \\
    \midrule
    non ideale & 1.0  & $\sim$84 & 0.41 & 100 & 0.34 & $\sim$99.5 & 0.19 \\
    non ideale & 0.5  & $\sim$87 & 0.39 & 100 & 0.25 & $\sim$98.5 & $\sim$0.17 \\ % TODO: verificare B-PDR forward da CSV (barra coperta dalla legenda)
    non ideale & 0.25 & $\sim$88 & 0.35 & 100 & 0.13 & $\sim$96   & 0.13 \\
    \bottomrule
  \end{tabular}
\end{table}

Tre osservazioni emergono dai dati. In primo luogo, per il single-hop
l'intervallo breve \`e un beneficio soltanto nel canale non ideale, dove la
maggiore frequenza di trasmissione offre pi\`u occasioni di superare i
collegamenti a bassa probabilit\`a ($D_{\%}$ dall'84 all'88\% passando da
1.0 a \SI{0.25}{\second}); nel canale ideale la scoperta \`e gi\`a satura al
proprio limite strutturale e l'unico effetto \`e il degrado del B-PDR (da
0.93 a 0.75). In secondo luogo, append paga l'intervallo aggressivo in modo
severo: a $BI=\SI{0.25}{\second}$ il collision rate nel canale ideale
raggiunge 0.77 a 100 boe (dal valore 0.27 con $BI=\SI{1.0}{\second}$) e il
B-PDR crolla a 0.23, senza alcun guadagno di scoperta, gi\`a completa
all'intervallo pi\`u lento. In terzo luogo, per forward l'intervallo breve
diventa controproducente anche sul piano della scoperta: a 100 boe $D_{\%}$
scende dal 100\% ($BI=\SI{1.0}{\second}$) al 98.5\% nel canale ideale e al
96\% circa in quello non ideale, poich\'e la congestione autoindotta
(collision rate fino a 0.82) sopprime proprio le ritrasmissioni che
dovrebbero diffondere l'informazione. Per entrambe le modalit\`a multi-hop
l'intervallo di \SI{1.0}{\second} risulta dunque il compromesso migliore:
la scoperta \`e gi\`a massima e ogni riduzione dell'intervallo si traduce
esclusivamente in contesa aggiuntiva.

% ---------------------------------------------------------------------
\section{Discussione}
\label{sec:discussione}
I risultati delineano un compromesso netto tra conoscenza della rete e
affidabilit\`a delle trasmissioni, il cui esito \`e per\`o asimmetrico tra le
due modalit\`a di inoltro. Il guadagno di scoperta \`e sostanzialmente
identico --- entrambe portano $D_{\%}$ al 100\% dove il single-hop si ferma
al 91\% (canale ideale) o all'84\% (canale non ideale) --- ma il prezzo
pagato \`e molto diverso. Append ottiene l'amplificazione informativa a
parit\`a di numero di trasmissioni, allungando i frame: il suo collision
rate cresce in modo contenuto (0.27 a 100 boe nel canale ideale) e il B-PDR
si mantiene su valori utilizzabili (0.74). Forward paga invece ogni hop con
una trasmissione aggiuntiva: il traffico cresce con la densit\`a fino a
saturare il canale (collision rate 0.74) e il B-PDR crolla a 0.26, un
terzo di quello di append a parit\`a di scoperta. La dominanza di append \`e
sistematica: a ogni densit\`a, intervallo e condizione di canale
considerati, append eguaglia la scoperta di forward con affidabilit\`a
superiore. L'unico vantaggio residuo di forward \`e nella fase transitoria a
bassissima contesa, dove le ritrasmissioni esplicite completano la scoperta
dei nodi pi\`u lontani senza dipendere dalla lunghezza delle liste
annunciate.

Il canale non ideale non ribalta il quadro, ma ne sposta i termini: il tetto
di affidabilit\`a imposto dal modello a due soglie (B-PDR $\le 0.45$ gi\`a in
assenza di contesa) ridimensiona il divario tra le modalit\`a, mentre la
ridondanza dell'inoltro si rivela un meccanismo di compensazione efficace,
mantenendo la scoperta sopra il 99\% dove il single-hop perde ulteriori
5--7 punti. Quanto all'intervallo di beaconing, la Sezione~\ref{sec:res-interval}
mostra che ridurlo amplifica il trade-off senza spostarne l'equilibrio: con
l'inoltro attivo la scoperta \`e gi\`a satura a $BI=\SI{1.0}{\second}$ e ogni
trasmissione aggiuntiva si converte in contesa, fino al caso limite di
forward la cui congestione autoindotta erode la scoperta stessa.

Restano infine i limiti della valutazione: la durata di \SI{120}{\second} e
l'area fissa di $500 \times 500\,\si{\metre\squared}$, il modello di canale a
due soglie privo di interferenze esterne e ostacoli, l'assenza dei protocolli
adattivi ADAB e ACAB dalla campagna qui presentata e il limite di 4 hop per
la modalit\`a forward, i cui effetti a diametri di rete maggiori non sono
stati esplorati.

\begin{table}[htbp]
  \centering
  \caption[Sintesi qualitativa del confronto tra modalit\`a]{Sintesi
  qualitativa del confronto tra le modalit\`a di inoltro.}
  \label{tab:summary}
  \begin{tabular}{lccc}
    \toprule
    & \textbf{Single-hop} & \textbf{Append} & \textbf{Forward} \\
    \midrule
    Scoperta della rete ($D_{\%}$) & parziale (76--91\%) & completa & completa \\
    Velocit\`a di scoperta & lenta, senza plateau & molto rapida ($<$\SI{6}{\second}) & rapida \\
    Affidabilit\`a (B-PDR) & massima & intermedia & bassa in reti dense \\
    Contesa (collision rate) & trascurabile & moderata & elevata, cresce con la densit\`a \\
    Robustezza al canale non ideale & scarsa (perde 5--7 p.p.) & alta & alta \\
    Overhead & nessuno & frame pi\`u lunghi & trasmissioni aggiuntive \\
    \bottomrule
  \end{tabular}
\end{table}

% ---------------------------------------------------------------------
\section{Sintesi del capitolo}
\label{sec:risultati-sintesi}
Il capitolo ha valutato le modalit\`a di inoltro multi-hop su una griglia di
densit\`a, intervalli di beaconing e condizioni di canale. Entrambe le
modalit\`a portano la scoperta della rete a valori prossimi al 100\%, contro
il 76--91\% del beaconing single-hop, e la mantengono anche in presenza di
perdite probabilistiche del canale. Il forwarding aggregato (append) si \`e
dimostrato la strategia dominante: ottiene la stessa conoscenza del
forwarding singolo con un collision rate da due a tre volte inferiore,
poich\'e affida la diffusione dell'informazione all'allungamento dei frame
anzich\'e a trasmissioni aggiuntive. Il forwarding singolo resta competitivo
solo in reti sparse, mentre nelle reti dense la sua contesa autoindotta ne
erode sia l'affidabilit\`a sia, agli intervalli pi\`u aggressivi, la scoperta
stessa. Le implicazioni di questi risultati e i possibili sviluppi sono
discussi nel capitolo conclusivo.


\end{document}
